\section{Non-Ergodic Mess3 Mixture}

\subsection{Construction}

We construct a training set of $K$ Mess3 processes with distinct parameters $\{(\alpha_k, x_k)\}_{k=1}^K$. Each training sequence is generated entirely by one component, chosen uniformly at random. The model observes only the token sequence---the component identity is latent.

\paragraph{Component selection.} We choose components with well-separated entropy rates $H_k$ and spectral gaps $\zeta_k = 1 - 3x_k$ to ensure meaningfully different dynamics:

\begin{table}[h]
\centering
\small
\begin{tabular}{lcccc}
\toprule
 & $\alpha$ & $x$ & $H_k$ (nats) & $\zeta_k$ \\
\midrule
C0 & 0.95 & 0.03 & 0.41 & 0.91 \\
C1 & 0.85 & 0.05 & 0.69 & 0.85 \\
C2 & 0.65 & 0.15 & 1.03 & 0.55 \\
\bottomrule
\end{tabular}
\caption{Mess3 components used in the non-ergodic mixture. $H_k$ is the entropy rate (theoretical minimum CE), $\zeta_k$ is the spectral gap controlling belief convergence rate.}
\label{tab:components}
\end{table}

C0 is highly predictable with slow mixing; C2 is nearly uniform with fast mixing. This ensures the components are distinguishable from their token statistics, while spanning a range of predictive difficulty.

\subsection{Relevance to Language Models}
\label{sec:relevance}

The non-ergodic structure mirrors key properties of natural language:
\begin{itemize}
    \item \textbf{In-context learning as source discrimination}: Each sequence begins with ambiguity about which regime is active. The model must use early tokens to identify the generating process---analogous to how LLMs adapt to document topic, style, or task from context.
    \item \textbf{Covariate shift}: Different components produce different token distributions, so the model faces a covariate shift across sequences. The model must be robust to this shift, adapting its predictions in-context rather than relying on a single fixed strategy.
    \item \textbf{Hierarchical structure}: Natural language has latent structure at multiple scales (topic, syntax, discourse). The component/state hierarchy in our setup is a minimal model of this: component identity is a slow variable (constant within a sequence), while hidden state is a fast variable (evolving at each token).
\end{itemize}

\subsection{Experimental Setup}

\paragraph{Model.} GPT-2 style decoder-only transformer implemented with TransformerLens~\cite{nanda2022transformerlens}. Phase 1 (single Mess3): 3 layers, $d_{\text{model}} = 64$, 4 heads. Phase 2 (non-ergodic mixture): 4 layers, $d_{\text{model}} = 128$, 4 heads. Pre-norm (LayerNorm), learned positional embeddings, ReLU activation. Vocabulary: $\{$BOS, 0, 1, 2$\}$ (4 tokens).

\paragraph{Training.} Next-token cross-entropy loss with Adam ($\text{lr} = 5 \times 10^{-4}$, no weight decay). Fresh batch of 8192 sequences generated on-the-fly at each step (no fixed dataset). Context length: 16 tokens + BOS. Phase 1: 10K steps. Phase 2: 20K steps.

\paragraph{Bayes-optimal baseline.} We compute the Bayes-optimal cross-entropy at each context position via Monte Carlo sampling: this is the CE achieved by a perfect Bayesian observer who maintains exact belief states. For the mixture, we compute both the \emph{oracle} baseline (knows the component identity) and the \emph{Bayesian} baseline (must infer it from observations).

\paragraph{Analysis methods.} (1) Cumulative explained variance (CEV) of residual stream activations via PCA. (2) Linear regression from activations to ground-truth belief states ($R^2$, RMSE). (3) Meta-belief regression (posterior over components). (4) Per-position and per-layer probing to track synchronization dynamics. (5) Attention pattern analysis.
